\section{Finiteness of the renormalized perturbation expansion}\label{sec:fin}

The renormalization constants
$Z$, $Z_3$ and $\tilde{Z}_3$ are now assumed to be such
that the singularities of the
correlation functions of the fields $(\ren{A})_{\mu},\ren{c}$
and $\ren{\cbar}$ at $D=4$
cancel to all orders of the renormalized coupling.
Our aim in this section is to show that
the correlation functions involving the bulk fields
at positive flow time are then finite too and thus
do not require further renormalization.

\subsection{Renormalized perturbation theory}

We first need to reorganize the perturbation expansion
in terms of the renormalized parameters and fields.
In this form,
the expansion is generated by an action $S_0+\Delta S$
of the renormalized fields,
which includes the counterterms $\Delta S$ required to
cancel the poles of the Feynman diagrams at $D=4$.

Explicitly $S_0$ and $\Delta S$ are obtained by
expressing the bare parameters and fields
in the total action
\begin{equation}
  \Stot=S+\Sgf+\Sgh+\Sfl+\Sflgh
  \label{eq:7.1}
\end{equation}
through the renormalized ones and by setting
\begin{align}
  S_0=\left.\Stot\right|_{Z=Z_3=\tilde{Z}_3=1},
  \label{eq:7.2}\\[2.5pt]
  \Delta S=\Stot-S_0+\Delta\Sbc,
  \label{eq:7.3}
\end{align}
where
\begin{align}
   \Delta\Sbc=2\int\rmd^Dx\,\tr\bigl\{
   (Z^{1/2}Z_3^{1/2}-1)L_{\mu}(0,x)(\Ar)_{\mu}(x)
   \notag\\[2pt]
   +
   (\tilde{Z}_3Z^{1/2}Z_3^{1/2}-1)\dbar(0,x)\ren{c}(x)\bigr\}.
   \label{eq:7.4}
\end{align}
The boundary conditions are then
\begin{equation}
  \left.B_{\mu}\right|_{t=0}=(\ren{A})_{\mu},
  \qquad
  \left.d\right|_{t=0}=\ren{c},
  \label{eq:7.5}
\end{equation}
and the Feynman rules derived from the action $S_0$
thus coincide with those discussed in sect.~\ref{sec:perturbation}
(apart from the fact that the bare parameters $g_0$ and $\lambda_0$
are replaced by $\mu^{\eps}g$ and $\lambda$, respectively).

Note that the boundary conditions \eqref{eq:7.5} differ from
the ones imposed on the bare fields (a product of
renormalization factors is missing).
The counterterm $\Delta\Sbc$ must be included
in the action of the renormalized fields
to correct for this.
It amounts to adding two-point vertices to the
Feynman rules, whose effect on the
gauge-field and ghost propagators is
equivalent to a change of the boundary
conditions by the missing renormalization factors.

\subsection{Absence of bulk counterterms}\label{ssec:nobulk}

In the renormalized perturbation expansion,
singularities at $D=4$ (if any) appear for the first time
at some loop order $l\geq1$.
Since the interactions are local, and since
all propagators are tempered distributions in position space
with singularities only at coinciding arguments,
the divergent parts of the correlation functions at
this loop order are expected to be
such that they can be canceled by local counterterms.
Similarly to the case of an ordinary
field theory on a half-space studied by Symanzik \cite{Symanzik},
the counterterms can be localized either in the bulk or at
the boundary of the half-space.

In the theory considered here,
bulk counterterms can be excluded from the outset.
In order to show this, first note that
the correlation functions of the bulk fields
$B_{\mu}$, $L_{\mu}$, $d$ and $\dbar$
are, at large flow times, given by
Feynman diagrams built from
flow lines
and flow vertices only.
All diagrams
of this kind are directed trees or products such trees,
each tree ending at one of the $B$ and $d$ fields in
the correlation function considered. Moreover, the
flow lines at the other ends of the trees must
start from the $L$ and $\dbar$ fields in the
correlation function.

Since there are no loop diagrams,
the correlation functions of the bulk fields are
non-singular at large flow times and do not require
renormalization.
Divergent bulk counterterms are therefore excluded.
Note, incidentally, that
correlation functions of local fields composed from $B$ and $d$
fields are finite too, because no loop
diagrams are generated when the arguments of some of
these fields coalesce.
The field $D_{\mu}d$ that appears in the BRS
identity \eqref{eq:6.23}, for example, is of this kind
and consequently does not need to be renormalized.

\subsection{Boundary counterterms}

The discussion in the previous
subsection shows that the structure of any divergent
parts of the correlation functions
must correspond to the insertion
of a counterterm localized at flow time zero.
Since the correlation functions of the renormalized fields
$(\ren{A})_{\mu},\ren{c}$ and $\ren{\cbar}$
are finite to all orders, divergences can only
arise from diagrams with at least one flow vertex
and thus at least one such vertex with an external flow line.
The possible counterterms are therefore proportional to
$L_{\mu}$ or $\dbar$.

Since the gauge coupling is dimensionless,
the counterterm
may not involve composite fields of dimension larger than $4$
(such terms would be irrelevant). Moreover, Lorentz symmetry,
global gauge invariance and the ghost number conservation
must be respected.
Since $L_{\mu}$ and $\dbar$ have dimension $3$
and all other fields dimension $1$, it follows that
the possible counterterms at $l$-loop order are of the form
\begin{equation}
   2g^{2l}\int\rmd^Dx\,\tr\bigl\{
   z_1L_{\mu}(0,x)(\Ar)_{\mu}(x)+
   z_2\dbar(0,x)\ren{c}(x)\bigr\}
   \label{eq:7.6}
\end{equation}
with some (singular) coefficients $z_1$ and $z_2$.
An $LB$ and a $\dbar d$ term should in principle be
included here, but in view of the boundary conditions \eqref{eq:7.5},
these terms are not independent and their inclusion
would be equivalent to a change of
the coefficients $z_1$ and $z_2$.

\subsection{Consequences of the BRS symmetry}

We now show that the boundary counterterm \eqref{eq:7.6} is excluded
by the BRS symmetry of the theory.
In terms of the renormalized fields and parameters, the
BRS identity \eqref{eq:6.23} reads
\begin{equation}
   \lambda\langle
   B^a_{\mu}(t,x)\partial_{\nu}(\ren{A})^b_{\nu}(y)
   \partial_{\rho}(\ren{A})^c_{\rho}(z)
   \rangle=
   -\langle(D_{\mu}d)^a(t,x)
   (\ren{\cbar})^b(y)\partial_{\rho}(\ren{A})^c_{\rho}(z)\rangle.
   \label{eq:7.7}
\end{equation}
The unusual asymmetric renormalization \eqref{eq:5.13} of the ghost fields
$c$ and $\cbar$ was chosen to ensure that the renormalization factors
drop out in this equation. From the point of view of the
$\SUn$ gauge theory, the standard and
the asymmetric renormalization of the ghost fields are equivalent,
but the situation
is different in the theory in $D+1$ dimensions, because only $c$ couples
to the bulk fields.

Equation \eqref{eq:7.7} holds as long as no extra counterterms need to be
added and thus up to loop order $l$ inclusive. Note that
all fields in the correlation functions in this
equation, including the composite field $D_{\mu}d$,
are renormalized fields that cannot be additionally
renormalized (cf.~sect.~\ref{ssec:nobulk}).
If the correlation functions
are singular at $l$-loop order, one must therefore
be able cancel the singularities by a counterterm of
the form \eqref{eq:7.6}.

The contribution of the counterterm \eqref{eq:7.6} to the correlation
functions is obtained by inserting the corresponding two-point
vertices in the tree diagrams in fig.~\ref{fig:brsdiag}. The insertions have
the effect of multiplying the values of the diagrams $1-5$ by
$g^{2l}$ times
\begin{equation}
   2z_1,\;z_1,\;z_1+z_2,\;z_1+z_2\;\hbox{and}\;z_2
   \label{eq:7.8}
\end{equation}
respectively. However,
recalling the values of the diagrams
quoted in appendix~\ref{app:C}, the sum of the diagrams
weighted by the factors \eqref{eq:7.8} turns out to violate
the BRS identity \eqref{eq:7.7} unless $z_1=z_2=0$.

The correlation functions in eq.~\eqref{eq:7.7} must therefore be finite
at $l$-loop order and all other (renormalized) correlation functions
must be non-singular too, because the addition to the action
of a counterterm of the form \eqref{eq:7.6} is excluded.
We have thus shown that the theory in $D+1$
dimensions does not require further renormalization.
